\begin{code}[hide]
module vqupit.sections.constructions where

open import Data.Sum using (_⊎_)
open import Level using (0ℓ)
open import Algebra.Bundles using (Group)
open import Word.Base using (Word ; WRel ; _ʰ' ; _ⁿ')
open import Presentation.Definitions using (_IsPresentationOf_)
open import Presentation.Construct.Base using ([_]ₗ ; [_]ᵣ ; ConjRelʷ ; EmptyRel)

infix 4 _≈₁_ _===₁_ _===₂_
infix 5 _⋄_⋄_

postulate
  proof-elided : ∀ {ℓ} {A : Set ℓ} → A
  A B N H : Set
  u v : Word A
  _===₁_ : WRel N
  _===₂_ : WRel H
  _≈₁_ : Word N → Word N → Set
  Γ Δ : Set
  conj : H → N → Word N
  G1⋊G2 : Group 0ℓ 0ℓ
  RelTwist : Set → Set → Set
  _∪_ : Set → Set → Set
\end{code}
\subsection{Products of presentations, and their presentation theorems}\label{sec:constructions}

The presentation just obtained is one factor of the groups
\S\ref{sec:composing} assembles; this last subsection recalls the
constructions that glue such factors together.  Relation sets over a
disjoint union of alphabets (\texttt{Presentation.Construct.Base}) are
built from a single primitive,
a three-way join whose \emph{mixed} component is the design parameter:

\begin{code}
data _⋄_⋄_ (Γ₁ : WRel A) (Γ₂ : WRel B) (Γ₃ : WRel (A ⊎ B)) : WRel (A ⊎ B) where
  left  : Γ₁ u v → (Γ₁ ⋄ Γ₂ ⋄ Γ₃) [ u ]ₗ [ v ]ₗ
  right : Γ₂ u v → (Γ₁ ⋄ Γ₂ ⋄ Γ₃) [ u ]ᵣ [ v ]ᵣ
  mid   : Γ₃ u v → (Γ₁ ⋄ Γ₂ ⋄ Γ₃) u v
\end{code}

Choosing the mixed relation recovers the classical constructions ---
free product ($\aid{EmptyRel}$), direct product ($\aid{CommRel}$:
generators of the two sides commute), semidirect product
($\aid{ConjRel}$/$\aid{ConjRelʷ}$: moving $h$ past $n$ conjugates $n$,
by a generator or by a word), and definitional extension
($\aid{SugarRel}$: a new generator desugars to a word) --- plus $n$-fold
iterations $\Gamma \mathrel{\aid{⊕\char94}} n$.  Congruence-lifting
lemmas (\aid{lefts}, \aid{rights}) embed the factors' congruences into
the join, and normal-form transports carry witnesses across these maps:
a normal form for one factor becomes a weak normal form for the union,
and a normal form for a presented monoid pulls back along any monoid
monomorphism into it.

For each construction, \texttt{Presentation.Construct.Properties} proves
two kinds of theorems.  \emph{Normal-form lifting}: given normal-form
witnesses for the factors $\Gamma$ and $\Delta$, with carriers
$\mathit{NF}_1$ and $\mathit{NF}_2$, the construction has a normal-form
witness whose carrier is $\mathit{NF}_1 \times \mathit{NF}_2$ for the
direct and semidirect products.  \emph{Presentation}: if the factors
present concrete groups, the constructed relation set presents the
constructed group.  With their premise telescopes elided, the direct,
semidirect and $n$-fold statements read:
\[
\begin{array}{l}
\aid{dpres} : (\Gamma \mathrel{\aidop{⋄}} \Delta \mathrel{\aidop{⋄}} \aid{CommRel})
   \ \aid{IsPresentationOf}\ (G_1 \times G_2) \\[3pt]
\aid{dpres} : (\Gamma \mathrel{\aidop{⋄}} \Delta \mathrel{\aidop{⋄}} \aid{ConjRelʷ}\;\mathit{conj})
   \ \aid{IsPresentationOf}\ (G_1 \rtimes G_2) \\[3pt]
\aid{presentation} : \forall n \to (\Gamma \mathrel{\aid{⊕\char94}} n)
   \ \aid{IsPresentationOf}\ (\otimes\aid{-group}\;n)
\end{array}
\]
where the semantic groups on the right are standard-library-style
bundles built by the library's \texttt{ForStdlib} layer.  The
semidirect theorem is the one \S\ref{sec:composing} uses, so we state it
in full.  Its mixed relation $\aid{ConjRelʷ}\;\mathit{conj}$ says that
moving a letter $h$ of the acting group past a letter $n$ of the normal
subgroup conjugates $n$, by a \emph{word}: $h \bullet n =
\mathit{conj}\;h\;n \bullet h$.  The theorem takes the two relation sets
and the action, two well-definedness conditions on the action --- it
respects the acting group's axioms in its first argument
(\aid{conj-hyph}) and the normal subgroup's in its second
(\aid{conj-hypn}) --- and presentations of the factors, and returns a
presentation of the semidirect product of the two \emph{semantic}
groups:

\begingroup\renewcommand{\AgdaCodeStyle}{\footnotesize}
\begin{code}
module _ (conj-hyph : ∀ {c d} n → c ===₂ d → (conj ʰ') c n ≈₁ (conj ʰ') d n)
         (conj-hypn : ∀ c {w v} → w ===₁ v → (conj ⁿ') c w ≈₁ (conj ⁿ') c v) where
  module Presentation (G1 G2 : Group 0ℓ 0ℓ)
    (p1 : Γ IsPresentationOf G1) (p2 : Δ IsPresentationOf G2) where
    dpres : (Γ ⋄ Δ ⋄ ConjRelʷ conj) IsPresentationOf G1⋊G2
\end{code}
\endgroup
\begin{code}[hide]
    dpres = proof-elided

extension-presentation : Set → Set → Set → Set → Set
\end{code}

\noindent Here \aid{G1⋊G2} is built by the library's standard-library
style semidirect product from an \aid{Action} record whose action is the
syntactic \aid{conj} transported through the two presentation
isomorphisms; its normal form is the product of the factors' normal
forms, and uniqueness reduces factor by factor to the injectivity of
the two factor interpretations (\aid{unfp₃}).  A further module proves the classical recipe for
presenting a group \emph{extension}: if $1 \to N \to G \to Q \to 1$ is a
short exact sequence, and presentations of $N$ and $Q$ are given together
with the data that realises it in $G$ --- how $G$ conjugates $N$, and the
correction in $N$ by which each lifted relator of $Q$ fails to hold ---
then $N$'s relations, word-valued conjugation rules, and the corrected
lifts of $Q$'s relators together present $G$ (cf.\ Proposition~2.55
of~\cite{holt2005handbook}).  A second mixed component, \aid{RelTwist},
twists each relator $u = v$ of the quotient by its correction word in
the normal subgroup, and the relation set of the recipe --- the one the
qupit paper invokes, in its pen-and-paper form, for its Section~4.2 ---
is
\begingroup\renewcommand{\AgdaCodeStyle}{\footnotesize}
\begin{code}
extension-presentation S R conj corr = S ⋄ EmptyRel ⋄ (ConjRelʷ conj ∪ RelTwist R corr)
\end{code}
\endgroup
and the library proves it as a presentation theorem for any group
extension realised by a short exact sequence and, in the form used in
\S\ref{sec:scalars}, for central extensions built from a cocycle.  Semidirect products are the special
case with trivial corrections.  These construction theorems are what
make the composition step of \S\ref{sec:composing} cheap: the semidirect
product of the Pauli and symplectic presentations and the central
extension by the scalars are \emph{assembled} from the factor
presentations, with no fresh coset enumeration at all.
